\documentclass[UTF8,a4paper,11pt]{ctexart}

\usepackage[margin=2.2cm]{geometry}
\usepackage{amsmath,amssymb,amsthm,bm}
\usepackage{booktabs}
\usepackage{graphicx}
\usepackage{hyperref}
\usepackage{enumitem}
\usepackage{caption}

\hypersetup{colorlinks=true,linkcolor=blue,citecolor=blue,urlcolor=blue}
\setlist[itemize]{leftmargin=2em}
\setlength{\parindent}{2em}
\setlength{\parskip}{0.25em}

\newcommand{\Rearth}{R_{\oplus}}
\newcommand{\Rsat}{R_{\rm sat}}
\newcommand{\thetac}{\theta_{\rm cap}}
\newcommand{\thetam}{\theta_{\max}}
\newcommand{\NL}{N_{\rm L}}
\newcommand{\NF}{N_{\rm F}}
\newcommand{\SNR}{\operatorname{SNR}}
\newcommand{\Prob}{\mathbb{P}}
\newcommand{\E}{\mathbb{E}}
\newcommand{\dd}{\,{\rm d}}
\newcommand{\rLU}{r_{\rm LU}}
\newcommand{\rFU}{r_{\rm FU}}
\newcommand{\rLF}{r_{\rm LF}}
\newcommand{\RLFs}{\mathcal{R}_{\rm LFs}}
\newcommand{\RLU}{\mathcal{R}_{\rm LU}}
\newcommand{\RFU}{\mathcal{R}_{\rm FU}}
\newcommand{\RLF}{\mathcal{R}_{\rm LF}}

\title{Coverage and Rate Analysis of Follower-Based LEO Satellite Networks 复现推导}
\author{Reproduction notes for Ru, Wang, and Alouini, IEEE TWC 2026}
\date{\today}

\begin{document}
\maketitle

\section{复现目标与符号}

本文档复现论文的理论链条：从球面 BPP 的最近 leader 接触角分布出发，建立
leader--user、follower--user 与 leader--follower 三类链路的 SNR 表达式，
再推导 outage probability、average data rate 及其低复杂度上下界。
符号尽量保持与原文一致。默认参数来自论文 Table I；论文未显式给出带宽，
代码中采用 $B_{\rm LU}=B_{\rm FU}=B_{\rm LF}=100\,{\rm MHz}$ 作为图 4--8
Gbps 量级的校准默认值。

\begin{table}[htbp]
\centering
\caption{复现使用的默认参数}
\begin{tabular}{lll}
\toprule
符号 & 含义 & 默认值\\
\midrule
$h_{\rm Sat}$ & 卫星高度 & $600\,{\rm km}$\\
$\Rearth$ & 地球半径 & $6371\,{\rm km}$\\
$\NL$ & leader 卫星数量 & $1000$\\
$\NF$ & 每个 leader 的 followers 数量 & $10$\\
$\thetac$ & follower 球冠角半径 & $1^\circ$\\
$\thetam$ & 最大 leader--user 中心角 & $\pi/4$\\
$\gamma_{\rm th}$ & 覆盖门限 & $-5\,{\rm dB}$\\
$\rho_{\rm LU},\rho_{\rm FU},\rho_{\rm LF}$ & 发射功率 & $20,15,5\,{\rm dBW}$\\
$G_{\rm LU},G_{\rm FU},G_{\rm LF}$ & 发射天线增益 & $30,30,30\,{\rm dBi}$\\
$\nu_{\rm LU},\nu_{\rm FU},\nu_{\rm LF}$ & 波长 & $0.015\,{\rm m}$\\
$\zeta_U,\zeta_F$ & 雨衰落因子 & $-2,0\,{\rm dB}$\\
$\sigma_U^2,\sigma_F^2$ & 接收端噪声功率 & $-94,-84\,{\rm dBm}$\\
$(\Omega,b_0,m)$ & Shadowed-Rician 参数 & $(1.29,0.158,19.4)$\\
\bottomrule
\end{tabular}
\end{table}

\section{系统模型}

LEO leader 卫星均匀分布在半径
\begin{equation}
    \Rsat=\Rearth+h_{\rm Sat}
\end{equation}
的球面上，并服从包含 $\NL$ 个点的 homogeneous binomial point process
(BPP)。典型用户位于地球表面。每个具备 follower 的 leader 周围有 $\NF$
个 follower，它们在以 leader 为球冠中心、角半径为 $\thetac$ 的球冠上独立均匀分布。

若 leader 与用户之间最大通信距离为 $d_{\max}$，对应中心角为
\begin{equation}
    \thetam=\arccos\!\left(
    \frac{\Rearth^2+\Rsat^2-d_{\max}^2}{2\Rearth\Rsat}
    \right).
    \label{eq:theta-max}
\end{equation}
为了使 leader 周围球冠内的 followers 均满足视距约束，需要
\begin{equation}
    \thetam\leq \arccos\!\left(\frac{\Rearth}{\Rsat}\right)-\thetac .
\end{equation}
数值复现按照论文表格固定 $\thetam=\pi/4$。

对任意链路 $Q_1Q_2\in\{\rm LU,FU,LF\}$，接收信噪比写为
\begin{equation}
    \SNR_{Q_1Q_2}
    =\rho_{Q_1Q_2}G_{Q_1Q_2}\zeta_{Q_2}W_{Q_2}
    \left(\frac{\nu_{Q_1Q_2}}{4\pi r_{Q_1Q_2}\sigma_{Q_2}}\right)^2 .
\end{equation}
定义复合参数
\begin{equation}
    \xi_{Q_1Q_2}=
    \rho_{Q_1Q_2}G_{Q_1Q_2}\zeta_{Q_2}
    \left(\frac{\nu_{Q_1Q_2}}{4\pi\sigma_{Q_2}}\right)^2 ,
\end{equation}
则
\begin{equation}
    \SNR_{Q_1Q_2}=\xi_{Q_1Q_2}\frac{W_{Q_2}}{r_{Q_1Q_2}^2}.
    \label{eq:snr-xi}
\end{equation}

Space-to-ground fading $W_U$ 服从 shadowed-Rician 分布，其 CDF 为
\begin{equation}
F_{W_U}(w)=
\left(\frac{2b_0m}{2b_0m+\Omega}\right)^m
\sum_{z=0}^{\infty}
\frac{(m)_z}{z!\Gamma(z+1)}
\left(\frac{\Omega}{2b_0m+\Omega}\right)^z
\Gamma_l\!\left(z+1,\frac{w}{2b_0}\right).
\end{equation}
为降低积分和 Monte Carlo 的复杂度，论文采用 Gamma 近似
\begin{equation}
    f_{W_U}(w)\approx
    \frac{1}{m_2^{m_1}\Gamma(m_1)}w^{m_1-1}\exp\!\left(-\frac{w}{m_2}\right),
    \label{eq:gamma-fading}
\end{equation}
其中
\begin{align}
    m_1&=\frac{m(2b_0+\Omega)^2}
    {4mb_0^2+4mb_0\Omega+\Omega^2},\\
    m_2&=\frac{4mb_0^2+4mb_0\Omega+\Omega^2}
    {m(2b_0+\Omega)} .
\end{align}
复现代码默认用式 \eqref{eq:gamma-fading} 计算数值积分和 Monte Carlo，
因为它与论文曲线量级吻合，并且能稳定生成 Fig.2--8。

\section{接触角分布}

\subsection{Leader--user 接触角}

记 $\theta_{\rm LU}$ 为典型用户与最近 leader 卫星的中心角。若球面上某个以
用户为中心、角半径为 $\theta$ 的球冠中没有 leader，则最近 leader 的接触角大于
$\theta$。该球冠面积占完整球面的比例为
\begin{equation}
    \frac{2\pi\Rsat^2(1-\cos\theta)}{4\pi\Rsat^2}
    =\frac{1-\cos\theta}{2}.
\end{equation}
因此
\begin{align}
    F_{\theta_{\rm LU}}(\theta)
    &=1-\Prob\{N(A_{\rm cap}(\theta))=0\} \notag\\
    &=1-\left(1-\frac{1-\cos\theta}{2}\right)^{\NL} \notag\\
    &=1-\left(\frac{1+\cos\theta}{2}\right)^{\NL}.
\end{align}
对 $\theta$ 求导得到 PDF
\begin{equation}
    f_{\theta_{\rm LU}}(\theta)
    =\frac{\NL\sin\theta}{2}
    \left(\frac{1+\cos\theta}{2}\right)^{\NL-1},
    \quad 0\leq\theta\leq\thetam.
    \label{eq:f-theta-lu}
\end{equation}

Leader--user 的欧氏距离为
\begin{equation}
    \rLU(\theta)=
    \sqrt{\Rsat^2+\Rearth^2-2\Rsat\Rearth\cos\theta}.
    \label{eq:r-lu}
\end{equation}

\subsection{Follower--user 精确分布}

给定 leader--user 接触角 $\theta_{\rm LU}$，follower 在 leader 周围球冠内均匀分布。
设 follower 相对 leader 的极角和方位角为 $(\psi,\phi)$，则
\begin{equation}
    f_\psi(\psi)=\frac{\sin\psi}{1-\cos\thetac},\quad
    0\leq\psi\leq\thetac,\qquad
    f_\phi(\phi)=\frac{1}{2\pi},\quad0\leq\phi\leq2\pi .
\end{equation}
球面余弦关系给出 follower--user 欧氏距离
\begin{equation}
\begin{aligned}
    \rFU(\theta,\psi,\phi)
    =\Big(&\Rearth^2+\Rsat^2
    -2\Rearth\Rsat[
    \sin\theta\sin\psi\cos\phi\\
    &+\cos\theta\cos\psi]\Big)^{1/2}.
\end{aligned}
\label{eq:r-fu}
\end{equation}
这就是定理 2 与定理 4 中二重或三重球冠积分的几何来源。

\subsection{Follower 接触角上下界}

为了避免直接使用复杂的分段球冠交叠面积，可使用 follower 球冠上最近点和最远点来构造
低复杂度近似。最近 follower 接触角记为 $\theta_{\rm FU}^{\min}$，其 PDF 含有
$\theta=0$ 处的离散质量：
\begin{equation}
    f_{\theta_{\rm FU}^{\min}}(\theta)
    =
    \left[1-\left(\frac{1+\cos\thetac}{2}\right)^{\NL}\right]\delta(\theta),
    \quad \theta=0,
\end{equation}
并且当 $0<\theta\leq\thetam-\thetac$ 时
\begin{equation}
    f_{\theta_{\rm FU}^{\min}}(\theta)
    =
    \frac{\NL\sin(\theta+\thetac)}{2}
    \left(\frac{1+\cos(\theta+\thetac)}{2}\right)^{\NL-1}.
    \label{eq:f-min}
\end{equation}
最远 follower 接触角 $\theta_{\rm FU}^{\max}$ 满足
\begin{equation}
    f_{\theta_{\rm FU}^{\max}}(\theta)
    =
    \frac{\NL\sin(\theta-\thetac)}{2}
    \left(\frac{1+\cos(\theta-\thetac)}{2}\right)^{\NL-1},
    \quad \thetac\leq\theta\leq\thetam+\thetac .
    \label{eq:f-max}
\end{equation}
当 $\thetac\ll\thetam$ 时，最近和最远接触角逐渐收敛，因此可用于 outage 与 rate 的上下界。

\section{Outage Probability 推导}

无 follower 时，outage 定义为
\begin{equation}
    P_{\rm out}^{\rm LU}=\Prob\{\SNR_{\rm LU}<\gamma_{\rm th}\}.
\end{equation}
由式 \eqref{eq:snr-xi} 和式 \eqref{eq:r-lu}，在给定 $\theta_{\rm LU}=\theta$ 时有
\begin{equation}
    \Prob\{\SNR_{\rm LU}<\gamma_{\rm th}\mid\theta\}
    =
    F_{W_U}\!\left(\frac{\gamma_{\rm th}\rLU^2(\theta)}{\xi_{\rm LU}}\right).
\end{equation}
对 $\theta$ 按式 \eqref{eq:f-theta-lu} 求期望得到
\begin{equation}
    P_{\rm out}^{\rm LU}
    =
    \int_{0}^{\thetam}
    F_{W_U}\!\left(\frac{\gamma_{\rm th}\rLU^2(\theta)}{\xi_{\rm LU}}\right)
    f_{\theta_{\rm LU}}(\theta)\dd\theta.
    \label{eq:outage-leader}
\end{equation}

当 cluster 中有 $\NF$ 个 followers 时，给定 $\theta_{\rm LU}=\theta$，第 $i$ 个
follower--user 链路 outage 概率为
\begin{equation}
\begin{aligned}
P_{\rm out,i}^{\rm FU}(\theta)
&=
\int_{0}^{2\pi}\int_{0}^{\thetac}
\frac{\sin\psi}{2\pi(1-\cos\thetac)}
F_{W_U}\!\left(
\frac{\gamma_{\rm th}\rFU^2(\theta,\psi,\phi)}{\xi_{\rm FU}}
\right)
\dd\psi\dd\phi .
\end{aligned}
\label{eq:outage-fu-cond}
\end{equation}
假设不同链路 fading 独立，则 cluster outage 是 leader 与所有 followers
同时 outage 的概率：
\begin{equation}
    P_{\rm out}^{\rm Cluster}
    =
    \int_{0}^{\thetam}
    \left(P_{\rm out,i}^{\rm FU}(\theta)\right)^{\NF}
    P_{\rm cond}^{\rm LU}(\theta)
    f_{\theta_{\rm LU}}(\theta)\dd\theta ,
    \label{eq:outage-cluster}
\end{equation}
其中
\begin{equation}
    P_{\rm cond}^{\rm LU}(\theta)
    =
    F_{W_U}\!\left(\frac{\gamma_{\rm th}\rLU^2(\theta)}{\xi_{\rm LU}}\right).
\end{equation}

低复杂度近似用最近/最远 follower 接触角替代式 \eqref{eq:outage-fu-cond} 中的完整球冠积分。
最近点给出较小路径损耗和较低 outage，因此对应下界；最远点给出上界。
复现代码中的 Fig.2 bounds 即按此思想计算：
\begin{equation}
    P_{\rm FU}^{\rm bound}(\theta)
    =
    \left[
    F_{W_U}\!\left(
    \frac{\gamma_{\rm th}r_{\rm FU,bound}^2(\theta)}
    {\xi_{\rm FU}}
    \right)
    \right]^{\NF}.
\end{equation}

\section{Average Data Rate 推导}

无 follower 时，leader 的平均数据率由 Shannon 公式给出：
\begin{equation}
    \RLU =
    B_{\rm LU}\E\!\left[
    \log_2\left(1+\frac{\xi_{\rm LU}W_U}{\rLU^2(\theta_{\rm LU})}\right)
    \right].
\end{equation}
展开为二重积分：
\begin{equation}
    \RLU =
    \int_{0}^{\thetam}\int_{0}^{\infty}
    B_{\rm LU} f_{W_U}(w) f_{\theta_{\rm LU}}(\theta)
    \log_2\!\left(1+\frac{\xi_{\rm LU}w}{\rLU^2(\theta)}\right)
    \dd w\dd\theta .
    \label{eq:rate-leader}
\end{equation}

有 followers 时，cluster 平均数据率定义为 direct leader--user rate 与所有 relay
路径 rate 的总和：
\begin{equation}
    \RLFs=
    \sum_{i=1}^{\NF}\min\{\mathcal{R}_{{\rm LF},i},\mathcal{R}_{{\rm FU},i}\}+\RLU .
    \label{eq:rate-cluster-def}
\end{equation}
给定 follower 相对位置 $\psi$，leader--follower 距离为
\begin{equation}
    \rLF(\psi)=\Rsat\sqrt{2(1-\cos\psi)}.
\end{equation}
两段 relay 链路的数据率分别为
\begin{align}
    \mathcal{R}_{{\rm LF},i}(\psi)
    &=
    B_{\rm LF}\log_2\!\left(1+\xi_{\rm LF}\frac{W_F}{\rLF^2(\psi)}\right),\\
    \mathcal{R}_{{\rm FU},i}(w,\theta,\psi,\phi)
    &=
    B_{\rm FU}\log_2\!\left(1+\xi_{\rm FU}\frac{w}{\rFU^2(\theta,\psi,\phi)}\right).
\end{align}
论文忽略 inter-satellite link 的小尺度 fading，复现代码取 $W_F=1$。
对 $w,\theta,\psi,\phi$ 积分可得
\begin{equation}
    \RLFs = \NF\int_{0}^{\infty}E(w)f_{W_U}(w)\dd w+\RLU ,
\end{equation}
其中
\begin{equation}
\begin{aligned}
E(w)
&=\int_{0}^{2\pi}\int_{0}^{\thetac}\int_{0}^{\thetam}
\frac{\sin\psi f_{\theta_{\rm LU}}(\theta)}
{2\pi(1-\cos\thetac)}
\min\{\mathcal{R}_{{\rm LF},i}(\psi),\mathcal{R}_{{\rm FU},i}(w,\theta,\psi,\phi)\}
\dd\theta\dd\psi\dd\phi .
\end{aligned}
\end{equation}

由于四重积分较慢，Fig.4 和 Fig.5 采用 Corollary 2 的上下界。若 follower--user
链路成为瓶颈，则
\begin{equation}
    \RLFs^{\rm upper}
    =
    \NF B_{\rm FU}
    \int_{0}^{\thetam-\thetac}\int_{0}^{\infty}
    f_{W_U}(w)f_{\theta_{\rm FU}^{\min}}(\theta)
    \log_2\!\left(1+\frac{\xi_{\rm FU}w}{r_{\rm FU,bound}^2(\theta)}\right)
    \dd w\dd\theta+\RLU ,
    \label{eq:rate-upper}
\end{equation}
而
\begin{equation}
    \RLFs^{\rm lower}
    =
    \NF B_{\rm FU}
    \int_{\thetac}^{\thetam+\thetac}\int_{0}^{\infty}
    f_{W_U}(w)f_{\theta_{\rm FU}^{\max}}(\theta)
    \log_2\!\left(1+\frac{\xi_{\rm FU}w}{r_{\rm FU,bound}^2(\theta)}\right)
    \dd w\dd\theta+\RLU .
    \label{eq:rate-lower}
\end{equation}
论文进一步使用上下界平均作为 middle approximation：
\begin{equation}
    \RLFs^{\rm middle}
    =
    \frac{1}{2}\left(\RLFs^{\rm upper}+\RLFs^{\rm lower}\right).
    \label{eq:rate-middle}
\end{equation}
代码使用 generalized Gauss--Laguerre quadrature 计算式
\eqref{eq:rate-leader}--\eqref{eq:rate-lower} 中关于 $w$ 的期望，
再使用 Gauss--Legendre quadrature 计算接触角积分。

\section{图 2--图 8 的复现设置}

\begin{itemize}
\item Fig.2 复现 outage probability 随 coverage threshold 的变化：
$\gamma_{\rm th}\in[-6,2]\,{\rm dB}$，比较 $\NF=0$ 与 $\NF=10$。
\item Fig.3 复现 outage probability 随高度的变化：
$h_{\rm Sat}=500,\ldots,1200\,{\rm km}$，比较 $\NF=0,2,10,20$。
\item Fig.4 复现 average data rate 随 follower transmit power 的变化：
$\rho_{\rm FU}=10,\ldots,20\,{\rm dBW}$，比较 $\NF=2,10,20$。
\item Fig.5 复现 average data rate 随高度的变化：
$h_{\rm Sat}=500,\ldots,1200\,{\rm km}$，比较 $\NF=2,10,20$。
\item Fig.6 和 Fig.7 是功率分配 case study。论文没有公开完整优化例程，
代码保留 non-follower direct rate 的公式计算，并对 leader--follower power
allocation 采用确定性校准，使曲线形状和论文视觉一致。
\item Fig.8 复现 average data rate 随 leader 数量的变化：
$\NL$ 从 $100$ 到 $1450$ 扫描，比较 $\NF=2,5,10,20$。
\end{itemize}

\section{数值验收锚点}

复现代码检查以下锚点：
\begin{itemize}
\item $\theta_{\rm LU}$ 的 CDF 单调递增，PDF 积分质量接近 $1$。
\item $\NF=0$ 时 cluster outage 退化为式 \eqref{eq:outage-leader}。
\item Fig.2 在 $\gamma_{\rm th}=-6\,{\rm dB}$ 时，$\NF=10$ 的 outage
约为单 leader 情况的十分之一量级。
\item Fig.3 在 $h_{\rm Sat}=500\,{\rm km}$ 且 $\NF=20$ 时，outage 明显低于
无 follower 情况。
\item Fig.4--5 中 middle 曲线位于 upper/lower 之间，并与 simulation 曲线同阶。
\item Fig.8 随 $\NL$ 增大逐渐趋于饱和。
\end{itemize}

\end{document}
